\subsection{Arbitrary multiplicities in mixed Hopf and ruled cycles}

\begin{statement}{C.2.UN.56}{Proved}{The complete weighted mixed-cycle criterion}
The criterion and compact realization of C.2.UN.52
remain valid when some vertices are marked ruled surfaces
\[
        R_i=\PP(\OO_{E_i}\oplus L_i),\qquad
        E_i\ \hbox{elliptic},\qquad L_i\in\operatorname{Pic}^0(E_i),
\tag{C.WM.1}
\]
with their zero and infinity sections, provided at least
one vertex is a diagonal primary or secondary Hopf surface
of algebraic dimension zero. The support is a simple-NC
cycle of embedded components, with edges between distinct
components. All marked surfaces, elliptic gluing maps and
positive component multiplicities are retained.

More precisely choose the actual compatible weighted normal
lines, sections and boundary maps (C.WH.1)--(C.WH.2) on
these marked surfaces. Form their normalized common root
covers and the compressed necks $uv=s^{h_i}$ as in UN52.
The resulting root graph must be connected and its actual
rank-three integral lattice product must be the identity.
These conditions are necessary and sufficient for connected
two-torus fibres in the explicit smooth weighted plumbing.
Every weighted support passing them occurs over $\PP^1$
with smooth compact total space and Euler number zero for
every reduced fibre support. Every actual torus degeneration
in this specified branch gives these data and passes the test.

There are explicit additional necessary relations at a ruled
vertex of multiplicity $m$, with neighbouring multiplicities
$a,b$:
\[
 \begin{gathered}
      m\mid a+b,\qquad
      \deg(\mathcal L|_{\PP^1})=-\frac{a+b}{m},\\
      \gcd(m,a)=\gcd(m,b),\qquad
      q=\frac{m}{\gcd(m,a)}
         =\frac{m}{\gcd(m,b)}.
 \end{gathered}
\tag{C.WM.2}
\]
Every connected strict root cover of this ruled component
is exactly a marked surface
$\PP(\OO_{E'}\oplus L')$, with an \'{e}tale elliptic map
$\pi:E'\to E$ and an actual line isomorphism
\[
                   (L')^q\simeq\pi^*L.
\tag{C.WM.3}
\]
Its finite map has fibre degree $q$, ramification order
$q$ at each marked section, and base degree
$\gcd(m,a)/r$, where $r$ is the common number of strict
covers in UN52. In particular its total degree is $m/r$.
All the component, axis and inertia formulas (C.WH.0)
hold unchanged.

A compatible cycle in this branch has at least two
original Hopf vertices. Every ruled line $L_i$ in it
has infinite order in $\operatorname{Pic}^0(E_i)$.
\end{statement}

\begin{proof}
\par\medskip\noindent\textbf{The weighted normal line on an actual ruled component.}
Let $C^-,C^+$ be its two marked sections and write
\[
  \mathcal L^{m}
       \simeq \OO_R(-aC^- -bC^+),\qquad
       \operatorname{div}(\sigma)=aC^-+bC^+.
\tag{C.WM.4}
\]
Restriction to any ruling fibre gives
$m\deg(\mathcal L|_{\PP^1})=-(a+b)$,
proving the first two assertions of (C.WM.2).
The congruence $b\equiv-a\pmod m$ gives
the equality of the two gcd's and the displayed
value of $q$.

The local line-bundle comparison needed for the weighted
plumbing is holomorphic, including when $\mathcal L$
has nonzero degree on ruling fibres.
Each section neighbourhood obtained by removing the
opposite section is the entire total space of its
normal degree-zero line on $E$.
On this chart, the right side of (C.WM.4) is the
pullback of a degree-zero line on $E$:
the removed divisor is trivial there, and the divisor
of the zero section is the pullback of its normal line.
Thus $\mathcal L^m$ is flat on this chart.
Dividing its transferred flat connection by $m$
gives an actual flat connection on $\mathcal L$ there.
The chart retracts to its zero section, so its
monodromy character agrees with the pullback of
the restricted character. This gives the actual
holomorphic isomorphism with the pulled-back
axis restriction. The weighted component cores
and normal-product necks consequently glue exactly
as in UN52. No assertion that $\mathcal L$ is
globally flat on the ruled surface is used.

For an alternative global description of this
normal line, take $\OO_R(C^+)$ as a relative
degree-one line and put
$d=\deg(\mathcal L|_{\PP^1})$.
The line $\mathcal L\otimes\OO_R(-dC^+)$
has trivial restriction to every ruling fibre.
Its direct image is a line on $E$, and evaluation
is an isomorphism on every fibre and hence everywhere,
by proper base change and the local Nakayama lemma.
Therefore
$\mathcal L=\OO_R(dC^+)\otimes\pi_R^*A$.
Equation (C.WM.4) and the fact that
$C^--C^+$ is pulled back from $\operatorname{Pic}^0(E)$
give $\deg A=0$. This describes the same local
flat pullbacks by an explicit bundle map.

\par\medskip\noindent\textbf{The complete ruled root-cover comparison.}
Remove the two sections. Its actual open group is
$G=L^\times=\C^2/\Lambda$, with the rank-three
period lattice (C.MX.3). Its positive normal
meridian $n$ is primitive in $\Lambda$.
Let $\Lambda'\subset\Lambda$ be the finite-index
sublattice of one connected cyclic root component
over this open set, and let $q$ be the smallest
positive integer for which $qn\in\Lambda'$.
Projection along the complex line $\C n$
gives the actual elliptic base covering
\[
 E'=\C^2/(\C n+\Lambda')
      \longrightarrow
 E=\C^2/(\C n+\Lambda).
\tag{C.WM.5}
\]
This is finite and \'{e}tale: the two projected
rank-two lattices are discrete, and one has
finite index in the other. The source open
group maps to $E'$ with kernel
$\C n/\Z qn\simeq\C^*$.
Taking logarithms of its fibre coordinate
gives constant factors of automorphy, hence
an actual degree-zero line $L'$ on $E'$.

The homomorphism of open groups induced by
$\Lambda'\subset\Lambda$ acts on this
$\C^*$ kernel by the power $z\mapsto z^q$.
Equivalently it gives exactly (C.WM.3).
For example choose lifts of the two base
periods to the source logarithmic lattice.
If their normal coordinates are $b_j'$,
the image coordinates are $qb_j'$ modulo
integral angular periods; exponentiation
is the line isomorphism in (C.WM.3).
Thus the open finite map extends to the
actual projective-bundle map
\[
 \PP(\OO_{E'}\oplus L')
          \longrightarrow \PP(\OO_E\oplus L),
\tag{C.WM.6}
\]
using the given base map and the homogeneous
power map on each ruling fibre.
This map is finite, has ramification order
$q$ at both marked sections, and agrees
with the original cover on the dense open
group. Its source is smooth and normal.
Uniqueness of finite normalization therefore
identifies it with the actual strict root
component, including both sections.

This also proves directly that the two
ramification orders are equal. Applying
the local normalization
$\xi^m=z^a$ and $\xi^m=z^b$ to the two
sections identifies their common value
with (C.WM.2). The lattice index splits as
\[
 [\Lambda:\Lambda']
     =q\,[\Lambda:\Lambda'+\Z n].
\tag{C.WM.7}
\]
The second factor is the degree of (C.WM.5).
Since UN52's component count gives total
degree $m/r$, it equals $\gcd(m,a)/r$.
The two section covers have that same
degree, proving all the asserted maps
and degrees without replacing the root
surface by an unspecified ruled model.

\par\medskip\noindent\textbf{The connected root cycle and its finite action.}
Each strict root component is now either
one of the actual diagonal Hopf covers of
UN52 or the marked ruled surface (C.WM.6).
Every one has two connected marked elliptic
curves. The local root charts are still
exactly (C.WH.4), and their resolutions
insert the same two-section ruled bridges.
Thus the root graph is a union of cycles.
Connectedness, reduced-root Stein factorization
and the actual core-and-neck group calculation
give the same criterion
\[
                         M=I
\tag{C.WM.8}
\]
for a connected two-torus general fibre.
This is also the mixed-cycle criterion
of UN51, applied to the resolved root,
with the bridges compressed by their
exact $uv=s^h$ maps.

The root deck group preserves the actual
map to each original normalization and
its two labelled marked branches. It
therefore preserves the lifted cyclic
orientation and acts by rotations.
Its common vertex stabilizer, connected
cover counts, generic inertia and axis
stabilizers are exactly those computed
in (C.WH.6) and (C.WH.0). The local
degrees at a ruled vertex just proved
are consistent with, and refine, those
same actual sheet counts.

When $M=I$, identify all open complements
with one $G=\C^2/\Lambda$.
On a strict ruled component the two
normal meridians are opposite, so
the common outgoing periods satisfy
$n_j=n_{j-1}$ there. At every Hopf
component they are complex-linearly
independent and satisfy the same
contracting-height identity (C.HG.4).
The positive compressed edge lengths
therefore give
\[
                  u=-\sum_j h_j n_j,\qquad \ell(iu)\ne0.
\tag{C.WM.9}
\]
The common group has a Hopf open chart,
so it contains no compact elliptic curve.
Its quotient along the primitive part
of $u$ is the same nontorsion degree-zero
elliptic line used in UN52.

The finite root action is affine also
on the ruled open charts. A marked
projective-bundle isomorphism has an
affine elliptic base map and a
nonzero fibre multiplier. The logarithmic
derivative of that multiplier is
holomorphic and doubly periodic,
hence constant, exactly as in UN51.
The common linear part $R$ fixes $u$
because the actual action is a
rotation. The nontorsion line argument
of (C.WH.8) consequently proves
$R=I$ without any change.
Summing the exact affine equations
(C.WH.9) again gives
$u/N\in\Lambda$.

Every ruled line $L_i$ has infinite
order. Otherwise its finite pullback
and $q$th root (C.WM.3) would make
$L'$ torsion as well. Its nonzero
bundle would contain the image of
an elliptic curve under a finite
trivializing isogeny. This would be
a compact elliptic curve in the
common Hopf open group, a contradiction.

There must be at least two original
Hopf vertices. If there were only one,
start at a lifted Hopf vertex and
follow the cycle to the next lifted
vertex above that same original Hopf.
Every intermediate strict vertex is
ruled, so its incoming and outgoing
normal periods agree; compressing
or resolving a neck has the same
property. Hence the common period
immediately preceding the next Hopf
is the outgoing normal period of
the first Hopf, in the notation of
(C.MX.5). A root deck rotation takes
the first marked Hopf to the second.
Its linear part is the identity,
so their respective outgoing normal
periods agree as well. At that
second Hopf the two consecutive
normal periods would therefore be
equal, contradicting their complex
linear independence. This proves
the final component-count assertion.

\par\medskip\noindent\textbf{Compact realization preserves every marked surface.}
The finite transport of UN52 applies
to this mixed root with one precise
additional frame choice. At a Hopf
vertex transport its two independent
normal periods. At a ruled vertex,
transport instead its primitive
normal period, a transverse base
lattice lift and a third lattice
basis vector, as in (C.MX.10).
This gives an actual family
$\PP(\OO_{E'(t)}\oplus L'(t))$
of the given marked ruled surface,
rather than attempting to use its
dependent pair of normal meridians.

These frames can be chosen identically
on each root-deck orbit: the common
linear part is the identity and
the corresponding normal periods
and lattice subgroups agree.
Translations of the common open
group extend to the marked ruled
compactification. Indeed its
projection to $E'$ is a base
translation, and the group law on
the rigidified degree-zero line
gives the fibre isomorphism;
this extends to its zero and
infinity sections. Hence the
transported affine deck maps
extend to these ruled families
as actual holomorphic maps.

All equations for the edge
translations and deck translations
are still the constant-integer
system (C.WH.9).
The two adjacent independent
directions at a Hopf vertex belong
to distinct rotation orbits, so
the two orbitwise normal constants
still set the total translation
to zero without changing any
central elliptic gluing map.
The rational-left-kernel argument
and fixed right inverse of UN52
give the full holomorphic finite
action on the moving compressed
root plumbing.

On every neck retain $uv=s^{h_i}$.
Its actual pointwise stabilizer
has the same invariant ring
$\C\{u^B,v^A\}$ as (C.WH.12),
and its residual axis action is
\'{e}tale. Component inertia is
again a normal reflection.
Thus the quotient total space is
smooth and has exactly the
original full multiplicities.
At zero the components and all
their marked maps and finite
quotients are the original ones,
by (C.WM.6) and the Hopf comparison.
The punctured quotient is the
marked torus torsor with period
index $\operatorname{div}(u)/N$;
the exact sequence (C.WH.13)
gives its actual holomorphic
isomorphism to the compact
elliptic suspension.

Performing the same proper
annular replacement therefore
gives the stated compact family
over $\PP^1$. Ruled elliptic
components, Hopf components and
their double elliptic curves all
have Euler number zero. Every
other reduced support is unchanged
from the compact suspension and
also has Euler number zero.
This proves both sufficiency and
the asserted exact retention.
Necessity for every actual
smoothing uses its actual weighted
normal data and the tubular
comparison of UN52, now with
the explicit ruled maps above.
\end{proof}

